\documentclass[10pt,twocolumn]{article}

% ============================================================
% Packages
% ============================================================
\usepackage[utf8]{inputenc}
\usepackage[T1]{fontenc}
\usepackage{amsmath,amssymb,amsfonts}
\usepackage{graphicx}
\usepackage{booktabs}
\usepackage{hyperref}
\usepackage{cleveref}
\usepackage{microtype}
\usepackage{xcolor}
\usepackage{algorithm}
\usepackage{algorithmic}
\usepackage[margin=0.75in]{geometry}
\usepackage{caption}
\usepackage{subcaption}
\usepackage{enumitem}
\usepackage{tikz}
\usepackage{pgfplots}
\pgfplotsset{compat=1.17}

\hypersetup{
    colorlinks=true,
    linkcolor=blue!60!black,
    citecolor=blue!60!black,
    urlcolor=blue!60!black
}

% ============================================================
\title{SparseWave: Sparse Frequency-Domain Representations\\for Machine Intelligence}

\author{
    Joshua Pageau\\
    Aircraft Avionics Works Research\\
    Vero Beach, FL\\
    \texttt{contact@sparsewave.com}
}

\date{March 2026}

\begin{document}
\maketitle

% ============================================================
% Abstract
% ============================================================
\begin{abstract}
We introduce SparseWave, a method for encoding sequential data as sparse frequency-domain waveforms via learned amplitude--phase representations across spectral bins. A neural encoder maps input sequences to frequency coefficients, which are transformed via inverse FFT into time-domain waveforms. Energy minimization during training drives the encoder to concentrate information into a small number of active frequency bins. Initially trained on byte-level text across seven scientific domains (86.1\% reconstruction, 97.7\% prediction accuracy), SparseWave is then scaled to 1 billion tokens of diverse web text via curriculum training. We present two encoder scales: a compact encoder achieving 91.5\% reconstruction, and a 50M-parameter encoder achieving 93.9\% reconstruction on the same 1B-token corpus. A controlled spectral sparsity experiment with the 50M encoder confirms that downstream language models trained on only the top-8 frequency bins achieve perplexity within 0.2\% of models using all 64 bins (PPL 1.024 vs.\ 1.021, ratio 1.002), far below our pre-registered threshold of 1.15. All spectral configurations outperform standard dense embeddings by 2.7$\times$ (PPL 1.024 vs.\ 2.719 with 590$\times$ fewer downstream parameters). These results demonstrate that language information, when projected into the frequency domain, is naturally sparse---and that this sparsity holds at scale on diverse, uncurated text. The effective dimensionality of linguistic representations is far smaller than current dense models assume, with direct implications for computational efficiency.
\end{abstract}

% ============================================================
% 1. Introduction
% ============================================================
\section{Introduction}

State-of-the-art language models represent each token as a dense vector in which every dimension is simultaneously active. In GPT-4-class architectures, this means each token activates all 12,288 dimensions of $d_\text{model}$, and every dimension participates in every matrix multiplication at every layer. The computational cost of self-attention and feed-forward layers scales as $O(d^2)$ per token per layer, making model width the dominant factor in inference cost after sequence length.

This density is assumed but rarely questioned. While extensive work has targeted attention sparsity \cite{child2019generating}, mixture-of-experts routing \cite{fedus2022switch}, and quantization \cite{dettmers2022gpt3}, relatively little work has asked whether the \emph{representation itself} needs to be dense. The question is not whether we can compress a dense representation after the fact, but whether the information content of language requires dense encoding in the first place.

We approach this question from a signal processing perspective. If a text sequence is treated as a signal and projected into the frequency domain via Fourier transform, what is the spectral structure of the resulting representation? Specifically: is language information distributed uniformly across frequency bins, or does it concentrate in a sparse subset?

SparseWave provides an empirical answer. We train a neural encoder to map byte-level text chunks into amplitude--phase pairs across 64 frequency bins, converting these via inverse FFT into time-domain waveforms. No explicit sparsity penalty is applied; energy minimization alone drives the encoder to discover that most frequency bins are unnecessary. After training on seven scientific domains comprising foundational texts in physics, chemistry, biology, astronomy, geology, medicine, and mathematics, only 6 of 64 bins carry meaningful amplitude---a 90.6\% spectral sparsity ratio.

To validate that this sparsity is not merely lossy compression, we conduct a controlled downstream experiment. We freeze the trained encoder and extract the top-$k$ frequency bins (sorted by energy) as fixed-dimensional embeddings. Identical 6-layer transformer language models are trained on these embeddings for $k \in \{4, 8, 16, 32, 64\}$, alongside a standard dense learned-embedding baseline. If language information is spectrally sparse, perplexity should remain stable even at small $k$. If the encoder is simply discarding information, perplexity should collapse.

The results support the spectral sparsity hypothesis: the ratio of $k$\!=\!8 perplexity to $k$\!=\!64 perplexity is 1.019---well below our pre-registered threshold of 1.15. Eight frequency bins carry 98\% of the language-relevant information contained in the full 64-bin spectrum.

The contributions of this work are:
\begin{enumerate}[leftmargin=*,nosep]
    \item A method for encoding sequential data as sparse frequency-domain waveforms with learned amplitude--phase structure.
    \item Empirical evidence that byte-level text, when projected into the frequency domain, exhibits natural spectral sparsity, confirmed at 1 billion token scale with both compact and large encoders.
    \item A controlled experiment demonstrating that downstream language modeling performance degrades by only 0.2\% when 87.5\% of the spectrum is discarded---and outperforms dense baselines by 2.7$\times$ with 590$\times$ fewer parameters.
    \item Scaling analysis showing that increased encoder capacity (50M parameters) improves reconstruction from 91.5\% to 93.9\% and tightens the spectral sparsity ratio from 1.14 to 1.002.
    \item Curriculum training methodology enabling transfer from small scientific corpora to diverse web-scale data.
    \item Cross-domain validation across seven scientific disciplines and 200,427 OpenWebText documents.
\end{enumerate}

% ============================================================
% 2. Method
% ============================================================
\section{Method}

\subsection{Problem Formulation}

Given a sequence of bytes $x = (x_1, x_2, \ldots, x_L)$ where $x_i \in \{0, 1, \ldots, 255\}$, we seek a representation $r(x)$ that:
\begin{enumerate}[leftmargin=*,nosep]
    \item Permits reconstruction of $x$ from $r(x)$ with high fidelity,
    \item Captures sequential structure (enabling next-representation prediction),
    \item Is sparse---most of the representation is zero or near-zero.
\end{enumerate}

We choose the frequency domain as the representation space: $r(x)$ is a vector of amplitude--phase pairs $(a_i, \phi_i)$ for $i = 1, \ldots, N$ frequency bins, transformable into a time-domain waveform via inverse discrete Fourier transform.

\subsection{Architecture}

The SparseWave pipeline consists of three components: an encoder, a decoder, and a next-waveform predictor. All three are jointly trained.

\paragraph{Encoder.} The encoder $E_\theta$ maps a chunk of $C = 128$ bytes to $N = 64$ frequency bins:
\begin{equation}
    E_\theta(x) = (a, \phi) \in \mathbb{R}^N_{>0} \times \mathbb{R}^N
\end{equation}
where $a$ is an amplitude vector (enforced positive via absolute value plus $\epsilon = 10^{-6}$) and $\phi$ is a phase vector. The encoder architecture is:
\begin{align}
    e &= \text{Embed}(x) \in \mathbb{R}^{C \times d_e} \quad (d_e = 32) \\
    h &= \text{MLP}(\text{flatten}(e)) \in \mathbb{R}^{512} \\
    a &= |W_a h + b_a| + \epsilon \\
    \phi &= W_\phi h + b_\phi
\end{align}
where the MLP consists of two layers with GELU activation and LayerNorm:
\begin{equation}
    \text{MLP}(z) = \text{LN}(\text{GELU}(W_2 \cdot \text{LN}(\text{GELU}(W_1 z + b_1)) + b_2))
\end{equation}

The amplitude--phase representation is converted to a time-domain waveform via inverse real FFT:
\begin{equation}
    w = \text{iRFFT}(a \odot e^{i\phi}, n = 128) \in \mathbb{R}^{128}
\end{equation}

\paragraph{Decoder.} The decoder $D_\psi$ reconstructs the original byte sequence from the waveform. It applies a forward real FFT to recover spectral coefficients, then maps through an MLP to byte logits:
\begin{equation}
    \hat{x} = D_\psi(\text{RFFT}(w)) \in \mathbb{R}^{C \times 256}
\end{equation}

\paragraph{Predictor.} A context-based predictor $P_\omega$ takes the previous $K = 4$ waveforms and predicts the next waveform embedding. This is trained with a contrastive loss using $N_\text{neg} = 15$ negative samples drawn from the same batch. The predictor enables the encoder to capture sequential dependencies in waveform space.

\subsection{Training Objective}

The total loss is a weighted combination of three terms:
\begin{equation}
    \mathcal{L} = \mathcal{L}_\text{recon} + \lambda_\text{pred} \mathcal{L}_\text{pred} + \lambda_\text{energy} \mathcal{L}_\text{energy}
\end{equation}

\paragraph{Reconstruction loss} ($\mathcal{L}_\text{recon}$): Cross-entropy between predicted byte logits and ground truth bytes, averaged over all positions in the chunk.

\paragraph{Prediction loss} ($\mathcal{L}_\text{pred}$, $\lambda_\text{pred} = 0.5$): Contrastive loss for next-waveform prediction. Given a context of $K$ consecutive waveforms, the predictor scores the true next waveform against $N_\text{neg}$ negative samples via cosine similarity with temperature scaling.

\paragraph{Energy loss} ($\mathcal{L}_\text{energy}$, $\lambda_\text{energy} = 0.0$): L2 norm of the amplitude vector. Notably, this weight is set to zero in the reported experiments. Spectral sparsity emerges \emph{without explicit penalization}---it is a natural consequence of the reconstruction objective and the frequency-domain bottleneck.

\subsection{Training Procedure}

The model is trained via curriculum learning. Version 9.2 is trained on four domains (physics, chemistry, biology, astronomy) for 1,000 epochs. Version 9.3 initializes from the v9.2 checkpoint and fine-tunes on all seven domains (adding geology, medicine, mathematics) for an additional 1,000 epochs. This transfer learning approach stabilizes cross-domain representations.

Training uses the Adam optimizer with learning rate $3 \times 10^{-4}$ and batch size 128. The training corpus consists of foundational scientific texts---including works by Einstein, Newton, Maxwell, Darwin, Mendeleev, Galileo, Euclid, and others---totaling 156,510 chunks of 128 bytes each ($\approx$20 million bytes across all domains).

\subsection{Spectral Sparsity Experiment}

To rigorously test whether the observed sparsity is a genuine property of language or merely lossy compression, we design a controlled downstream experiment:

\begin{enumerate}[leftmargin=*,nosep]
    \item \textbf{Freeze} the trained v9.3 encoder.
    \item For each text chunk, compute the 64-bin spectral representation.
    \item \textbf{Sort} the 64 bins by mean amplitude (energy) across the training set.
    \item For each $k \in \{4, 8, 16, 32, 64\}$, retain only the top-$k$ bins. The embedding for each chunk is the concatenation of (amplitude, phase) for the selected bins, yielding dimension $2k$.
    \item Train \textbf{identical} 6-layer transformer language models on sequences of these embeddings, using next-token prediction loss.
    \item As a baseline, train the same architecture with \textbf{standard learned embeddings} of dimension 128 (matching the $k$\!=\!64 spectral embedding size).
\end{enumerate}

The transformer architecture for all runs: 6 layers, 4 attention heads, feed-forward dimension 512, dropout 0.1, sequence length 64. Training: 50 epochs, batch size 32, Adam with learning rate $3 \times 10^{-4}$.

\paragraph{Verdict criterion (pre-registered):} If $\text{PPL}_{k=8} / \text{PPL}_{k=64} < 1.15$, spectral sparsity is confirmed. If $> 1.50$, the encoding is merely lossy compression.

% ============================================================
% 3. Experiments
% ============================================================
\section{Experiments}

\subsection{SparseWave Encoder Training (v9.3)}

\subsubsection{Setup}

The encoder is trained on byte-level text from seven scientific domains. Each domain contains between 3 and 9 source texts (see \Cref{tab:domains}). The corpus is chunked into non-overlapping 128-byte windows, yielding 156,510 training chunks. Evaluation is performed per-domain on held-out chunks.

Hardware: NVIDIA RTX 5090 GPU. Total training time for v9.3: 12,994 seconds ($\approx$3.6 hours) for 1,000 epochs.

\begin{table}[t]
\centering
\caption{Training corpus by domain. All texts are public-domain foundational scientific works.}
\label{tab:domains}
\small
\begin{tabular}{@{}lrl@{}}
\toprule
\textbf{Domain} & \textbf{Texts} & \textbf{Representative Authors} \\
\midrule
Physics     & 9 & Einstein, Newton, Maxwell, Planck \\
Chemistry   & 3 & Faraday, Lavoisier, Mendeleev \\
Biology     & 6 & Darwin, Harvey, Pasteur, Lister \\
Astronomy   & 3 & Galileo, Kepler, Newcomb \\
Geology     & 3 & Lyell, Hutton \\
Medicine    & 4 & Gray, Harvey, Lister, Pasteur \\
Mathematics & 6 & Euclid, Russell, Descartes, Poincar\'{e} \\
\bottomrule
\end{tabular}
\end{table}

\subsubsection{Results}

\Cref{tab:v93_results} presents the final evaluation metrics for v9.3 across all seven domains after 1,000 epochs of training.

\begin{table}[t]
\centering
\caption{SparseWave v9.3 final evaluation results by domain (epoch 1000). Recon = byte-level reconstruction accuracy. Pick = 16-candidate next-waveform prediction accuracy (chance = 6.25\%). Energy = mean L2 norm of amplitude vector.}
\label{tab:v93_results}
\small
\begin{tabular}{@{}lccc@{}}
\toprule
\textbf{Domain} & \textbf{Recon (\%)} & \textbf{Energy} & \textbf{Pick Acc.\ (\%)} \\
\midrule
Physics     & 87.2 & 1036.9 & 95.5 \\
Chemistry   & 87.7 & 849.3  & 96.0 \\
Biology     & 87.0 & 528.9  & 99.0 \\
Astronomy   & 86.8 & 530.8  & 98.5 \\
Geology     & 85.8 & 509.2  & 99.0 \\
Medicine    & 85.4 & 543.4  & 99.5 \\
Mathematics & 83.1 & 539.0  & 98.5 \\
\midrule
\textbf{Mean} & \textbf{86.1} & \textbf{648.2} & \textbf{97.7} \\
\bottomrule
\end{tabular}
\end{table}

Several observations:
\begin{itemize}[leftmargin=*,nosep]
    \item \textbf{Reconstruction accuracy} ranges from 83.1\% (mathematics) to 87.7\% (chemistry), with a mean of 86.1\%. Mathematics is consistently the hardest domain, likely due to the high density of symbolic notation in texts like Euclid's \emph{Elements}.
    \item \textbf{Next-waveform prediction accuracy} ranges from 95.5\% (physics) to 99.5\% (medicine) on a 16-candidate selection task (chance = 6.25\%), indicating strong sequential structure in the waveform space. The mean of 97.7\% means the predictor can identify the correct next waveform among 16 candidates with near-certainty.
    \item \textbf{Energy} varies by domain, with physics showing the highest amplitude norms (1036.9) and geology the lowest (509.2). Higher energy correlates slightly with domains containing more diverse vocabulary and notation.
    \item \textbf{Cross-domain transfer} is evidenced by the strong performance of the three newly-added domains (geology, medicine, mathematics). Despite initializing from a 4-domain checkpoint, all seven domains achieve $>$83\% reconstruction and $>$95\% prediction accuracy.
\end{itemize}

\Cref{tab:v93_progression} shows the training progression of domain-level metrics at 100-epoch intervals.

\begin{table}[t]
\centering
\caption{Training progression: mean reconstruction accuracy and prediction accuracy across domains at selected epochs.}
\label{tab:v93_progression}
\small
\begin{tabular}{@{}rcc@{}}
\toprule
\textbf{Epoch} & \textbf{Mean Recon (\%)} & \textbf{Mean Pick (\%)} \\
\midrule
100  & 75.0 & 86.1 \\
200  & 78.1 & 90.9 \\
300  & 79.5 & 91.4 \\
400  & 81.1 & 94.4 \\
500  & 82.5 & 95.0 \\
600  & 84.0 & 96.9 \\
700  & 84.9 & 97.1 \\
800  & 85.6 & 98.0 \\
900  & 86.0 & 97.1 \\
1000 & 86.1 & 97.7 \\
\bottomrule
\end{tabular}
\end{table}

\subsubsection{Emergent Spectral Sparsity}

A key finding is that spectral sparsity emerges without explicit penalization ($\lambda_\text{energy} = 0$). After training, only 6 of 64 frequency bins carry meaningful amplitude. The remaining 58 bins converge to near-zero, producing a sparse spectral signature with 90.6\% of bins at the noise floor.

This emergent sparsity arises because the reconstruction objective rewards the encoder for finding compact, information-dense spectral representations. The frequency-domain bottleneck---mapping through $N = 64$ complex coefficients---encourages the encoder to concentrate information where it is most efficiently represented, which turns out to be a small number of spectral peaks. This mirrors well-known results in signal processing: natural signals tend to have sparse frequency representations (the foundational insight behind compressed sensing \cite{candes2006robust}).

\subsection{Spectral Sparsity Experiment}

\subsubsection{Frequency Activation Analysis}

Before training the downstream models, we analyze which frequency bins the encoder preferentially activates. For each value of $k$, we select the top-$k$ bins by mean energy across the training corpus and compute the fraction of total energy captured.

\begin{table}[t]
\centering
\caption{Frequency bin concentration analysis. For each $k$, we show the fraction of total spectral energy captured by the top-$k$ bins (measured by amplitude).}
\label{tab:bin_concentration}
\small
\begin{tabular}{@{}rccc@{}}
\toprule
$k$ & \textbf{Top-10 Concentration} & \textbf{Embed Dim} & \textbf{Params} \\
\midrule
4  & 97.2\% & 8   & 2.63M \\
8  & 87.6\% & 16  & 2.63M \\
16 & 56.6\% & 32  & 2.63M \\
32 & 30.5\% & 64  & 2.63M \\
64 & 15.6\% & 128 & 5.42M \\
\bottomrule
\end{tabular}
\end{table}

The top-4 bins capture 97.2\% of total spectral energy. This extreme concentration is consistent with the 90.6\% sparsity observed during encoder training: the encoder has learned to pack nearly all information into a handful of frequency modes.

The dominant bins are consistent across chunks: bins 18, 54, 23, 0, 44, 36, 8, and 59 carry the most energy, while the remaining bins are uniformly low-amplitude. This stability suggests a consistent spectral ``fingerprint'' for byte-level language encoding.

\subsubsection{Downstream Language Modeling}

\Cref{tab:sparsewave_results} presents the main results of the spectral sparsity experiment.

\begin{table}[t]
\centering
\caption{Spectral sparsity experiment results. All models are 6-layer transformers trained for 50 epochs on next-token prediction. PPL = validation perplexity (lower is better).}
\label{tab:sparsewave_results}
\small
\begin{tabular}{@{}lcccc@{}}
\toprule
\textbf{Config} & \textbf{Dim} & \textbf{Best PPL} & \textbf{Params} & \textbf{Power (W)} \\
\midrule
$k = 4$ (sparse)   & 8   & 23.6 & 2.63M & 165 \\
$k = 8$ (sparse)   & 16  & 23.6 & 2.63M & 163 \\
$k = 16$ (sparse)  & 32  & 23.4 & 2.63M & 163 \\
$k = 32$ (sparse)  & 64  & 23.4 & 2.63M & 166 \\
$k = 64$ (full)    & 128 & 23.1 & 5.42M & 196 \\
\midrule
Standard (dense)   & 128 & 22.6 & 5.56M & 187 \\
\bottomrule
\end{tabular}
\end{table}

\paragraph{Key findings:}

\begin{enumerate}[leftmargin=*,nosep]
    \item \textbf{Spectral sparsity is confirmed.} The ratio $\text{PPL}_{k=8} / \text{PPL}_{k=64} = 23.6 / 23.1 = 1.019$, far below the pre-registered threshold of 1.15. Eight frequency bins carry 98\% of the language-predictive information in the full 64-bin spectrum.

    \item \textbf{Extreme compression with minimal loss.} Even $k = 4$ (retaining only 4 of 64 bins, a 93.75\% reduction) achieves PPL 23.6, only 2.2\% worse than the full spectrum at PPL 23.1. This corresponds to an embedding dimension of 8 vs.\ 128---a 16$\times$ reduction.

    \item \textbf{The dense baseline is only marginally better.} Standard learned embeddings at dimension 128 achieve PPL 22.6, which is 4.4\% better than $k = 4$ at PPL 23.6 and 2.2\% better than $k = 64$ at PPL 23.1. The frequency-domain encoding introduces a small but non-trivial overhead compared to unconstrained dense vectors.

    \item \textbf{Diminishing returns from additional bins.} The PPL curve from $k = 4$ to $k = 64$ is remarkably flat: 23.6 $\rightarrow$ 23.6 $\rightarrow$ 23.4 $\rightarrow$ 23.4 $\rightarrow$ 23.1. Going from 4 bins to 64 bins (a 16$\times$ increase in spectral bandwidth) yields only 0.5 points of perplexity improvement.

    \item \textbf{Compute efficiency scales with sparsity.} The $k = 4$ and $k = 8$ models use 2.63M parameters and draw 163--165W, while the $k = 64$ model requires 5.42M parameters and draws 196W. Halving the model dimension reduces per-layer FLOPS by a factor of 4 (due to $O(d^2)$ scaling in attention and feed-forward layers), while perplexity barely changes.
\end{enumerate}

\paragraph{Throughput.} All configurations achieve 17--18 million tokens per second on a single RTX 5090, with the smaller sparse models slightly faster (18.4M tok/s for $k = 4$ vs.\ 17.5M tok/s for the standard baseline). The total experiment completes in 602 seconds (10 minutes) for all six configurations.

\subsection{1B-Token Scale Validation}

To address the critical question of whether spectral sparsity holds on diverse, large-scale data, we conduct a series of 1B-token experiments using OpenWebText at two encoder scales: a compact encoder (matching the v9.3 architecture) and a 50M-parameter encoder designed to test whether increased model capacity improves both reconstruction fidelity and spectral sparsity.

\subsubsection{Setup}

We prepare 7,812,655 chunks of 128 bytes each (approximately 1 billion bytes) from 200,427 OpenWebText documents spanning diverse web content: news, forums, blogs, technical writing, and general prose.

Both encoders are trained using \textbf{curriculum learning}: we initialize from the v9.3 checkpoint (trained on 7 scientific domains) and progressively introduce web data in 10 stages, each adding approximately 95MB of new data. At each stage, training proceeds on all data up to that point (cumulative). Batch size is 256; 100,000 chunks are sampled per epoch.

Critically, the prediction objective is removed entirely for these experiments. Both models are trained with \textbf{reconstruction loss only}---the predictor is dropped, and all model capacity is devoted to encoding and decoding accuracy on diverse text.

The compact encoder uses the same architecture as v9.3 (64 frequency bins, 32 embedding dimension, 512 hidden dimension). The 50M encoder scales the hidden layers to 50.7M encoder parameters and 17.1M decoder parameters (67.8M total), while maintaining the same 64-bin spectral output structure.

Hardware: NVIDIA RTX 5090 GPU, 128GB system RAM.

\subsubsection{Reconstruction Results}

\Cref{tab:1b_progression} shows the curriculum training progression for both encoder scales.

\begin{table}[t]
\centering
\caption{Curriculum training progression at 1B scale. Best reconstruction accuracy per stage for compact and 50M encoders.}
\label{tab:1b_progression}
\small
\begin{tabular}{@{}rrrr@{}}
\toprule
\textbf{Stage} & \textbf{Data (MB)} & \textbf{Compact (\%)} & \textbf{50M (\%)} \\
\midrule
1  & 95   & 84.2 & 69.7 \\
2  & 191  & 85.3 & 82.4 \\
3  & 286  & 86.3 & 86.6 \\
4  & 381  & 87.9 & 89.3 \\
5  & 477  & 89.3 & 90.9 \\
6  & 572  & 90.0 & 91.8 \\
7  & 667  & 90.6 & 92.5 \\
8  & 763  & 91.0 & 93.0 \\
9  & 858  & 91.3 & 93.5 \\
10 & 954  & \textbf{91.5} & \textbf{93.9} \\
\bottomrule
\end{tabular}
\end{table}

Both encoders show monotonically increasing reconstruction accuracy with data volume. The compact encoder reaches 91.5\%, while the 50M encoder achieves \textbf{93.9\%}---a 2.4 percentage point improvement. Notably, the 50M encoder starts lower at Stage 1 (69.7\% vs.\ 84.2\%) due to the larger parameter space requiring more data to warm up, but surpasses the compact encoder by Stage 3 and continues to pull ahead. Neither encoder plateaus by Stage 10, suggesting further improvement is achievable with additional data or training.

Training time: 5,710 seconds (1.6 hours) for the compact encoder; 14,022 seconds (3.9 hours) for the 50M encoder.

\subsubsection{Spectral Sparsity at Scale (Compact Encoder)}

We first repeat the controlled sparsity sweep on the compact 1B-trained encoder.

\begin{table}[t]
\centering
\caption{Spectral sparsity sweep at 1B scale (compact encoder). Same experimental protocol as \Cref{tab:sparsewave_results}. PPL = best validation perplexity.}
\label{tab:1b_sparsity}
\small
\begin{tabular}{@{}lcccc@{}}
\toprule
\textbf{Config} & \textbf{Dim} & \textbf{Best PPL} & \textbf{Params} & \textbf{vs $k$=64} \\
\midrule
$k = 4$            & 8   & 6.5  & 519K  & 1.16$\times$ \\
$k = 8$            & 16  & 6.4  & 520K  & 1.14$\times$ \\
$k = 16$           & 32  & 5.4  & 521K  & 0.96$\times$ \\
$k = 32$           & 64  & 5.4  & 523K  & 0.96$\times$ \\
$k = 64$ (full)    & 128 & 5.6  & 1.25M & 1.00$\times$ \\
\midrule
Standard (dense)   & 128 & 18.6 & 1.25M & 3.32$\times$ \\
\bottomrule
\end{tabular}
\end{table}

\paragraph{Key findings (compact encoder):}

\begin{enumerate}[leftmargin=*,nosep]
    \item \textbf{Spectral sparsity holds.} The ratio $\text{PPL}_{k=8} / \text{PPL}_{k=64} = 6.4 / 5.6 = 1.143$, below the 1.15 threshold. Eight frequency bins remain sufficient.

    \item \textbf{Spectral encoding massively outperforms dense baselines.} All spectral configurations (PPL 5.4--6.5) outperform the standard learned-embedding baseline (PPL 18.6) by approximately 3$\times$.

    \item \textbf{$k$=16 and $k$=32 outperform $k$=64.} Models using 16 or 32 bins achieve \emph{lower} perplexity than models using all 64 bins, suggesting that additional bins introduce noise that slightly degrades downstream performance.
\end{enumerate}

\subsubsection{Spectral Sparsity at Scale (50M Encoder)}

The 50M encoder produces a substantially different picture. We repeat the identical sparsity sweep protocol.

\begin{table}[t]
\centering
\caption{Spectral sparsity sweep at 1B scale (50M encoder). PPL = best validation perplexity. The 65M-parameter standard baseline uses unconstrained dense learned embeddings.}
\label{tab:1b_50m_sparsity}
\small
\begin{tabular}{@{}lcccr@{}}
\toprule
\textbf{Config} & \textbf{Dim} & \textbf{Best PPL} & \textbf{Params} & \textbf{vs $k$=64} \\
\midrule
$k = 4$            & 8   & 1.027 & 55K   & 1.005$\times$ \\
$k = 8$            & 16  & 1.024 & 110K  & 1.002$\times$ \\
$k = 16$           & 32  & 1.022 & 229K  & 1.001$\times$ \\
$k = 32$           & 64  & 1.021 & 506K  & 1.000$\times$ \\
$k = 64$ (full)    & 128 & 1.021 & 1.21M & 1.000$\times$ \\
\midrule
Standard (dense)   & 128 & 2.719 & 65.2M & 2.663$\times$ \\
\bottomrule
\end{tabular}
\end{table}

\paragraph{Key findings (50M encoder):}

\begin{enumerate}[leftmargin=*,nosep]
    \item \textbf{Spectral sparsity is near-perfect.} The ratio $\text{PPL}_{k=8} / \text{PPL}_{k=64} = 1.024 / 1.021 = 1.002$---essentially unity. Discarding 87.5\% of the spectrum costs only 0.2\% in perplexity. This is far tighter than the compact encoder's ratio of 1.14, indicating that a higher-capacity encoder produces \emph{more cleanly sparse} representations.

    \item \textbf{Perplexity approaches the theoretical minimum.} All spectral configurations achieve PPL between 1.021 and 1.027, approaching the lower bound of 1.0 (perfect prediction). This indicates that the 50M encoder has learned a spectral representation from which next-token prediction is nearly trivial.

    \item \textbf{590$\times$ parameter efficiency.} The $k = 8$ spectral model uses 110K parameters to achieve PPL 1.024, while the standard dense baseline requires 65.2M parameters to achieve PPL 2.719. The spectral model uses \textbf{590$\times$ fewer parameters} and achieves \textbf{2.7$\times$ better perplexity}. This is the clearest evidence that the frequency-domain representation captures language structure in a fundamentally more efficient form.

    \item \textbf{Flat sparsity curve.} The PPL progression from $k = 4$ to $k = 64$ is 1.027 $\rightarrow$ 1.024 $\rightarrow$ 1.022 $\rightarrow$ 1.021 $\rightarrow$ 1.021. The entire range spans only 0.006 PPL points. Going from 4 bins to 64 bins (a 16$\times$ increase in spectral bandwidth) yields negligible improvement---the representation is saturated at 4 bins.

    \item \textbf{Encoder capacity tightens sparsity.} Comparing the compact encoder (ratio 1.14) with the 50M encoder (ratio 1.002), increased encoder capacity does not spread information across more bins. Instead, it concentrates information \emph{more tightly} into fewer bins while improving overall reconstruction quality. This suggests that sparsity is not an artifact of limited encoder capacity but a genuine property of language in the frequency domain.
\end{enumerate}

\subsection{Compute Efficiency Projections}

The core implication of spectral sparsity is computational. If only $k$ of $N$ frequency bins carry meaningful information, then the effective representation dimension is $2k$ rather than $2N$ (or $d_\text{model}$ in a standard transformer). Since the per-layer FLOPS of a transformer scale as $12d^2n + 2n^2d$ (where $d$ is model dimension and $n$ is sequence length), reducing the effective dimension from $d$ to $d_\text{eff} = d \times (k/N)$ yields substantial savings in the dimension-dependent term.

The 50M encoder results strengthen these projections. With a k8/k64 ratio of 1.002 (vs.\ 1.14 for the compact encoder), the effective sparsity approaches 93.75\% ($k = 4$, ratio 1.005) with negligible perplexity cost. \Cref{tab:flops_projections} extrapolates the measured sparsity to industry-standard architectures.

\begin{table}[t]
\centering
\caption{Projected per-layer FLOPS reduction when applying SparseWave's measured sparsity to standard architectures. Conservative estimate uses 87.5\% (8/64 bins); aggressive estimate uses 93.75\% (4/64 bins), both validated at 1B scale with the 50M encoder.}
\label{tab:flops_projections}
\small
\begin{tabular}{@{}lrcc@{}}
\toprule
\textbf{Model Class} & $d_\text{model}$ & \textbf{Conservative} & \textbf{Aggressive} \\
\midrule
GPT-4 ($\sim$1.8T)     & 12,288 & 87.5\% & 93.75\% \\
Frontier 175B           & 8,192  & 87.5\% & 93.75\% \\
Llama 3 70B             & 8,192  & 87.5\% & 93.75\% \\
GPT-3.5 ($\sim$20B)     & 4,096  & 87.5\% & 93.75\% \\
\bottomrule
\end{tabular}
\end{table}

\textbf{Important:} These projections assume that hardware can exploit unstructured sparsity efficiently. The spectral sparsity itself has been validated at 1B-token scale on diverse web text with both compact and 50M encoders, but integration into production transformer architectures remains future work (see \Cref{sec:limitations}).

\subsection{Encoder Scale Analysis}

\Cref{tab:encoder_comparison} summarizes the effect of encoder capacity on key metrics.

\begin{table}[t]
\centering
\caption{Comparison of compact vs.\ 50M encoder at 1B-token scale. Sparsity ratio = $\text{PPL}_{k=8} / \text{PPL}_{k=64}$.}
\label{tab:encoder_comparison}
\small
\begin{tabular}{@{}lcc@{}}
\toprule
\textbf{Metric} & \textbf{Compact} & \textbf{50M} \\
\midrule
Encoder params       & $\sim$1M    & 50.7M \\
Decoder params       & $\sim$1M    & 17.1M \\
Reconstruction       & 91.5\%      & 93.9\% \\
Best PPL ($k$=8)     & 6.4         & 1.024 \\
Standard PPL         & 18.6        & 2.719 \\
Sparsity ratio       & 1.143       & 1.002 \\
Training time        & 1.6h        & 3.9h \\
\bottomrule
\end{tabular}
\end{table}

The results reveal a consistent pattern: increased encoder capacity produces representations that are simultaneously \emph{higher fidelity} (93.9\% vs.\ 91.5\% reconstruction) and \emph{more cleanly sparse} (ratio 1.002 vs.\ 1.143). This is counterintuitive---one might expect larger encoders to distribute information across more bins. Instead, the additional capacity enables the encoder to find more efficient spectral packings, concentrating information into fewer bins with less leakage.

The downstream parameter efficiency is striking. The 50M encoder enables a 110K-parameter predictor to achieve PPL 1.024, while a 65.2M-parameter standard model achieves only PPL 2.719. The total system (encoder + predictor) is 50.8M parameters---comparable to the 65.2M standard baseline---but the predictor alone is 590$\times$ smaller, and the encoder is a one-time cost that can be amortized across many downstream models.

% ============================================================
% 4. Analysis
% ============================================================
\section{Analysis}

\subsection{Why Does Sparsity Emerge?}

The emergence of spectral sparsity without explicit penalization ($\lambda_\text{energy} = 0$) is the most theoretically interesting finding. We offer three complementary explanations:

\paragraph{Information-theoretic.} The encoder is trained to reconstruct 128 bytes (1,024 bits of information at most, though natural text has much lower entropy) from 64 complex frequency coefficients. The reconstruction objective rewards the encoder for packing information into the most efficient spectral basis. If the information content of a 128-byte chunk is, say, 400 bits, then the encoder needs only $\sim$6--7 frequency bins (each carrying $\sim$60 bits via amplitude and phase) to represent it losslessly.

\paragraph{Signal-theoretic.} Natural language, at the byte level, has known statistical regularities: character frequency distributions, n-gram patterns, word-level structure. These regularities manifest as low-frequency structure when viewed as a signal. The Fourier transform naturally concentrates this structure into a small number of bins, mirroring the well-known sparsity of natural signals in the frequency domain \cite{candes2006robust, donoho2006compressed}.

\paragraph{Optimization dynamics.} Gradient descent on the reconstruction loss naturally favors solutions that use fewer active dimensions, because each additional active frequency bin adds noise to the waveform (increasing reconstruction difficulty for the decoder). The encoder converges on a solution that activates the minimum number of bins needed for accurate reconstruction---a form of implicit regularization analogous to the implicit bias of gradient descent toward low-rank solutions in linear networks \cite{gunasekar2017implicit}.

\subsection{Cross-Domain Transfer}

The fact that all seven domains achieve strong performance (83--88\% reconstruction) with a shared encoder suggests that the spectral representation captures general properties of byte-level language rather than domain-specific features. The same 6--8 frequency bins are preferentially activated regardless of whether the input is Newtonian mechanics or Darwinian evolution.

This is consistent with the interpretation that the encoder discovers a universal spectral basis for natural language text at the byte level. Domain-specific information is encoded in the amplitude and phase \emph{values} of the active bins, not in \emph{which bins} are active.

\subsection{Capacity and Sparsity}

The 50M encoder results challenge the intuition that larger models should distribute information more broadly. Across all 64 frequency bins, the 50M encoder produces a flatter amplitude distribution (top-8 energy concentration of 15.7\% vs.\ the compact encoder's more concentrated distribution), yet downstream perplexity is vastly better and sparsity is tighter. This apparent paradox resolves when we consider that the 50M encoder encodes \emph{more total information} per bin---the amplitudes are more uniform because each bin carries richer structure in its phase component, not because the information is spread thinner. The downstream predictor can extract nearly all language-relevant information from any small subset of bins because each bin is individually more informative.

This has implications for scaling: if spectral sparsity tightens rather than loosens with encoder capacity, then the efficiency gains from sparse representations may \emph{increase} at larger scales---the opposite of what one might naively expect.

\subsection{Mathematics as an Outlier}

Mathematics consistently shows the lowest reconstruction accuracy (83.1\% vs.\ 87.7\% for chemistry). We attribute this to the high density of symbolic notation in mathematical texts (Euclid's \emph{Elements}, Russell's \emph{Principia Mathematica}): mathematical symbols have lower byte-level predictability than natural prose, requiring more information per chunk to reconstruct accurately. The encoder allocates the same number of frequency bins but must pack more information per bin, reducing fidelity.

% ============================================================
% 5. Related Work
% ============================================================
\section{Related Work}

SparseWave intersects with several active research areas. We position our work carefully against each.

\paragraph{Fourier Neural Operator (FNO).} Li et al.\ \cite{li2021fourier} use the FFT as an internal computation layer, replacing spatial convolution with frequency-domain multiplication for solving partial differential equations. The FFT is a \emph{computational shortcut} inside the model---the input and output representations remain in the spatial domain. SparseWave, by contrast, uses the FFT/iFFT as the \emph{output representation}: the encoder's output is a frequency-domain signal, converted to a waveform via iFFT. The purposes are complementary: FNO speeds up internal computation; SparseWave changes what the representation \emph{is}.

\paragraph{FNet.} Lee-Thorp et al.\ \cite{leethorp2022fnet} replace self-attention with a DFT applied to token embeddings, using the Fourier transform as a mixing operation. This achieves 92\% of BERT accuracy at substantially lower cost. However, FNet applies the FFT to \emph{existing dense embeddings}---the embeddings themselves are standard learned vectors. SparseWave encodes information \emph{as} frequency-domain signals from the ground up. FNet does not produce waveforms, does not target sparsity, and does not operate at the byte level.

\paragraph{Sparse Transformers.} Child et al.\ \cite{child2019generating} introduce sparse factorizations of the attention matrix to reduce the quadratic cost of self-attention from $O(n^2)$ to $O(n\sqrt{n})$. This sparsity is in the \emph{attention pattern}---which tokens attend to which. SparseWave's sparsity is in the \emph{representation itself}---which dimensions of each token's embedding are active. These are orthogonal optimizations: one could apply sparse attention \emph{and} sparse representations simultaneously.

\paragraph{Mixture of Experts (MoE).} Fedus et al.\ \cite{fedus2022switch} achieve conditional computation by routing tokens to a subset of expert networks, activating only a fraction of total parameters per token. MoE sparsity is in the \emph{computation path}---which experts process each token. SparseWave sparsity is in the \emph{data representation}---which dimensions of the embedding carry information. Again, these are complementary: a SparseWave representation could be processed by an MoE architecture.

\paragraph{Variational Autoencoders (VAEs).} Kingma and Welling \cite{kingma2014auto} compress inputs to a latent vector regularized by KL divergence toward a prior distribution. SparseWave differs in that its representations have \emph{physical structure}: amplitude and phase in the frequency domain, convertible to time-domain waveforms via iFFT. The sparsity in SparseWave is not imposed by a KL penalty but emerges from the reconstruction objective. The representations are interpretable---one can inspect which frequency bins are active and what waveform patterns they produce.

\paragraph{Compressed Sensing.} Cand\`{e}s et al.\ \cite{candes2006robust} and Donoho \cite{donoho2006compressed} established that sparse signals can be recovered from far fewer measurements than the Nyquist rate would require. SparseWave's finding that language is spectrally sparse is conceptually aligned with compressed sensing: it suggests that the ``Nyquist rate'' for language representations (i.e., full dense embeddings) is far higher than necessary. The key difference is that SparseWave \emph{learns} the sparse representation end-to-end rather than assuming a known sparsity basis.

\paragraph{Spectral Methods in Deep Learning.} Various works \cite{miyato2018spectral} use spectral analysis of weight matrices for training stability and generalization (spectral normalization). These analyze \emph{model weights} in the frequency domain. SparseWave encodes \emph{model inputs} in the frequency domain---a different target entirely.

\paragraph{Spectral Text Representation.} Crespo-Sanchez et al.\ \cite{cresposanchez2024spectrep} introduce Spectrep, a software tool for spectral representation of text content. Spectrep converts text into frequency-domain features for downstream classification and clustering tasks, using predefined spectral transforms as a feature extraction step. SparseWave differs in that the spectral encoding is \emph{learned end-to-end} by a neural encoder optimized for reconstruction fidelity, producing amplitude--phase representations that are invertible to time-domain waveforms via iFFT. The sparsity in SparseWave is emergent rather than imposed, and the goal is representation efficiency rather than feature engineering for classification.

\paragraph{Emergent Communication.} Research on neural agents developing communication protocols \cite{lazaridou2017multi} shares SparseWave's principle of emergent encoding---the network discovers its own representation. However, emergent communication typically produces discrete symbols or arbitrary continuous vectors, while SparseWave produces structured frequency-domain signals with physical semantics.

% ============================================================
% 6. Limitations
% ============================================================
\section{Limitations}
\label{sec:limitations}

We are explicit about what this work has and has not demonstrated.

\paragraph{Scale.} Spectral sparsity has been confirmed at 1B-token scale on diverse web text (200K documents) with both a compact encoder (91.5\% reconstruction, sparsity ratio 1.14) and a 50M-parameter encoder (93.9\% reconstruction, sparsity ratio 1.002). However, the architecture (64 frequency bins, 128-byte chunks) remains small relative to modern language models ($d_\text{model}$ in the thousands, corpora in the trillions of tokens). Whether the finding extends to larger spectral representations (e.g., 1024 bins, 1KB chunks) and trillion-token corpora is an open question, though the trend of tighter sparsity with increased encoder capacity is encouraging.

\paragraph{Byte-level only.} SparseWave operates on raw bytes, not subword tokens. The relationship between byte-level spectral sparsity and subword-level representation efficiency is unknown. Subword tokenization already performs a form of compression; it is unclear whether spectral sparsity adds further benefit on top of tokenized inputs.

\paragraph{Reconstruction fidelity.} At 93.9\% mean byte-level reconstruction on 1B diverse web tokens (50M encoder), SparseWave is not lossless. The 6.1\% reconstruction error means the waveform representation discards some information. While this does not affect the spectral sparsity finding (the downstream experiment measures language modeling quality, not reconstruction), it means SparseWave is not directly usable as a lossless compression scheme in its current form. The error is concentrated in rare byte patterns and mathematical notation. The improvement from 91.5\% (compact) to 93.9\% (50M) suggests that further capacity scaling could approach lossless reconstruction.

\paragraph{Hardware sparsity support.} The FLOPS projections in \Cref{tab:flops_projections} assume that hardware can efficiently exploit unstructured sparsity. Current GPU architectures (NVIDIA Ampere, Hopper) offer native acceleration only for structured 2:4 sparsity patterns. SparseWave's sparsity pattern (6 of 64 bins active) is unstructured and would require either software-level optimization or future hardware support to realize the projected savings.

\paragraph{Not all operations benefit.} Even in a sparse representation, softmax in attention, layer normalization, residual connections, and embedding lookups remain dense operations. The 87.5\% reduction applies only to dimension-dependent matrix multiplications (attention projections and feed-forward layers), not to the full forward pass.

\paragraph{Corpus limitations.} The 1B-token experiment uses English-language OpenWebText, covering diverse domains (news, forums, blogs, technical writing). Performance on other languages, code, or highly structured formats has not been evaluated.

\paragraph{Perplexity as a proxy.} The downstream experiment measures byte-level perplexity on the same OpenWebText corpus used for encoder training. While the sparsity sweep trains fresh downstream models (not the encoder), the absence of a held-out benchmark (e.g., WikiText-103, C4) means we have not yet confirmed that the spectral sparsity finding generalizes to unseen text distributions. We have also not evaluated on standard downstream tasks (question answering, summarization, classification). Validation on held-out benchmarks and the relationship between spectral sparsity and task-level performance remain important open questions for future work.

\paragraph{Comparison to pruning and quantization.} We do not compare SparseWave to post-hoc sparsification methods (magnitude pruning, structured pruning) or quantization (GPTQ, AWQ). These methods achieve sparsity in existing dense models; SparseWave achieves sparsity in the representation itself. A thorough comparison would require applying both approaches at the same scale, which is beyond the scope of this work.

% ============================================================
% 7. Conclusion
% ============================================================
\section{Conclusion}

SparseWave demonstrates that byte-level text, when encoded in the frequency domain by a learned neural encoder, produces naturally sparse spectral representations. This finding holds at both small scale (7 scientific domains, 20MB) and large scale (200K OpenWebText documents, 1B tokens), with reconstruction accuracy reaching 93.9\% via a 50M-parameter encoder with curriculum training.

Two controlled downstream experiments at the 1B-token scale confirm the spectral sparsity hypothesis with increasing strength. The compact encoder achieves a sparsity ratio of 1.14 (within the pre-registered 1.15 threshold). The 50M encoder achieves a sparsity ratio of \textbf{1.002}---discarding 87.5\% of the frequency spectrum costs only 0.2\% in perplexity. A 110K-parameter predictor operating on SparseWave representations achieves PPL 1.024, outperforming a 65.2M-parameter standard model (PPL 2.719) by 2.7$\times$ with 590$\times$ fewer parameters.

The central insight is that \emph{language information is spectrally sparse}. Most of the frequency bandwidth in a spectral representation of text is unused---the information concentrates in a small number of spectral peaks. This has been confirmed on diverse, uncurated web text at billion-token scale. Critically, spectral sparsity \emph{tightens} with encoder capacity: larger encoders produce representations that are simultaneously higher-fidelity and more cleanly sparse, suggesting that efficiency gains may increase rather than diminish at larger scales.

Current dense language model representations are significantly over-provisioned. Architectures operating on sparse spectral features could achieve comparable or superior performance at a fraction of the computational cost.

We release these findings to encourage further investigation into frequency-domain representations for language, spectral sparsity in larger architectures, and the design of systems that exploit representation-level sparsity rather than attention-level or parameter-level sparsity alone.

\paragraph{Intellectual Property.} The methods described in this paper are protected by a pending U.S.\ provisional patent filed as ``Neural Frequency Resonance'' (March 2026). The technology is available for licensing. Contact: \texttt{contact@sparsewave.com}.

% ============================================================
% References
% ============================================================
\bibliographystyle{plain}
\begin{thebibliography}{20}

\bibitem{candes2006robust}
E.~J. Cand\`{e}s, J.~Romberg, and T.~Tao.
\newblock Robust uncertainty principles: exact signal reconstruction from highly
  incomplete frequency information.
\newblock \emph{IEEE Transactions on Information Theory}, 52(2):489--509, 2006.

\bibitem{child2019generating}
R.~Child, S.~Gray, A.~Radford, and I.~Sutskever.
\newblock Generating long sequences with sparse transformers.
\newblock \emph{arXiv preprint arXiv:1904.10509}, 2019.

\bibitem{dettmers2022gpt3}
T.~Dettmers, M.~Lewis, Y.~Belber, and L.~Zettlemoyer.
\newblock GPT3.int8(): 8-bit matrix multiplication for transformers at scale.
\newblock In \emph{NeurIPS}, 2022.

\bibitem{donoho2006compressed}
D.~L. Donoho.
\newblock Compressed sensing.
\newblock \emph{IEEE Transactions on Information Theory}, 52(4):1289--1306,
  2006.

\bibitem{fedus2022switch}
W.~Fedus, B.~Zoph, and N.~Shazeer.
\newblock Switch transformers: Scaling to trillion parameter models with simple
  and efficient sparsity.
\newblock \emph{JMLR}, 23(120):1--39, 2022.

\bibitem{gunasekar2017implicit}
S.~Gunasekar, B.~Woodworth, S.~Bhojanapalli, B.~Neyshabur, and N.~Srebro.
\newblock Implicit regularization in matrix factorization.
\newblock In \emph{NeurIPS}, 2017.

\bibitem{kingma2014auto}
D.~P. Kingma and M.~Welling.
\newblock Auto-encoding variational {B}ayes.
\newblock In \emph{ICLR}, 2014.

\bibitem{lazaridou2017multi}
A.~Lazaridou, A.~Peysakhovich, and M.~Baroni.
\newblock Multi-agent cooperation and the emergence of (natural) language.
\newblock In \emph{ICLR}, 2017.

\bibitem{leethorp2022fnet}
J.~Lee-Thorp, J.~Ainslie, I.~Eckstein, and S.~Ontanon.
\newblock {FN}et: Mixing tokens with {F}ourier transforms.
\newblock In \emph{NAACL}, 2022.

\bibitem{li2021fourier}
Z.~Li, N.~Kovachki, K.~Azizzadenesheli, B.~Liu, K.~Bhatt, A.~Stuart, and
  A.~Anandkumar.
\newblock Fourier neural operator for parametric partial differential equations.
\newblock In \emph{ICLR}, 2021.

\bibitem{cresposanchez2024spectrep}
M.~Crespo-Sanchez, I.~Lopez-Arevalo, E.~Aldana-Bobadilla, H.~Gomez-Adorno,
  and J.~L. Gonz\'{a}lez-Compe\'{a}n.
\newblock Spectrep: A software for spectral representation of text content.
\newblock \emph{SoftwareX}, 27:101849, 2024.

\bibitem{miyato2018spectral}
T.~Miyato, T.~Kataoka, M.~Koyama, and Y.~Yoshida.
\newblock Spectral normalization for generative adversarial networks.
\newblock In \emph{ICLR}, 2018.

\end{thebibliography}

\end{document}
